\documentclass[11pt]{article}
\usepackage[T1]{fontenc}
\usepackage{lmodern}
\usepackage[margin=1.12in]{geometry}
\usepackage{microtype}
\usepackage{amsmath,amssymb,amsthm,mathtools}
\usepackage[hidelinks]{hyperref}
\newtheorem{theorem}{Theorem}
\providecommand{\tightlist}{\setlength{\itemsep}{0pt}\setlength{\parskip}{0pt}}
\setlength{\emergencystretch}{2em}
\allowdisplaybreaks[1]
\title{Positive metric entropy for billiards near a fixed ellipse}
\author{Kaiwen Jin\thanks{\href{mailto:yui.ui.siki@gmail.com}{\texttt{yui.ui.siki@gmail.com}}}}
\date{}
\hypersetup{pdftitle={Positive metric entropy for billiards near a fixed ellipse},pdfauthor={Kaiwen Jin}}
\begin{document}
\maketitle

\begin{abstract}
We construct smooth strictly convex planar billiards with positive
Kolmogorov--Sinai entropy for their invariant collision probability,
whose boundaries converge in \(C^\infty\) to the ellipse with semiaxes
\(2\) and \(\sqrt{3}\). The focal separatrix of this ellipse provides long
saddle passages joined by reflections on disjoint boundary arcs.
Changing these arcs realizes return maps close to the Berger--Turaev
stochastic-island model. We calibrate the periodic actions and then
restore the four bounding heteroclinic circles by varying actual
reflection profiles. The resulting positive-area islands give positive
entropy for the collision maps.
\end{abstract}

\section{Introduction}

In a planar convex billiard, straight flights alternate with specular
reflections at the wall. The collision map acts on boundary arclength and
outgoing angle. Its invariant area is the natural measure for asking
whether a set of initial conditions of positive probability is chaotic.
Elliptical tables are integrable: their confocal caustics organize the
collision dynamics, and their focal trajectories form a separatrix.

Smooth changes of the wall can produce hyperbolic basic sets. In particular,
Bessa, Del Magno, Lopes Dias, Gaiv\~ao and Torres proved positive topological
entropy for a \(C^2\)-open dense set of smooth convex bodies
\cite[Theorem 1.1 and Corollary 1.2]{Bessa}. Positive entropy for the
invariant collision probability calls for a positive-measure mechanism.
Among general smooth area-preserving surface maps, Berger and Turaev
obtained positive metric entropy by constructing and preserving a
stochastic island \cite[Theorem A]{BT}. Their perturbations act on the
phase map. Berger's Problem 2.17 asks for positive metric entropy when
the perturbation is a change of an elliptical billiard table
\cite[\S2.3]{BergerICM}.

\begin{theorem}\label{thm:main}
Let \(\gamma_E(\theta)=(2\cos\theta,\sqrt{3}\sin\theta)\) for
\(\theta\in S^1=\mathbb R/(2\pi\mathbb Z)\).
There exist \(C^\infty\) embeddings \(\gamma_j:S^1\to\mathbb R^2\) whose
images bound strictly convex domains \(D_j\), such that
\(\gamma_j\to\gamma_E\) in \(C^\infty\) and
\[
 h_{\mu_j}(B_j)>0\qquad\text{for every }j.
\]
Here \(B_j\) is the specular collision map of \(D_j\). With \(s\) the
arclength on \(\partial D_j\) and \(\varphi\in(0,\pi)\) the outgoing angle
from the positive tangent, its invariant probability is
\[
 d\mu_j=\frac{\sin\varphi}{2|\partial D_j|}\,ds\,d\varphi .
\]
\end{theorem}

The billiard near the focal separatrix makes repeated passes through
an unaltered saddle neighborhood. Between passes it reflects once at a
designated small arc of the wall. In rescaled coordinates this last
reflection supplies a shear, while the intervening passage supplies
the other part of a H\'enon factor. Physically disjoint arcs allow us
to compose such factors with a single wall. We first choose the finite
word and then make its arcs small. This order gives a sequence of
returns approaching the fixed stochastic-island model.

Approximating the model leaves its heteroclinic branches free to split.
To restore them, we use periodic chord lengths to remove the mean
defect. A boundary profile affects every visit to its arc with the same
scalar function, so the linearized equation must include all such
visits. Unfolding the first
visits lets us align the same profile with the native phase. Later visits
on the longer branch give a contracting contribution in a derivative
norm. Fixed-loss tame estimates let us solve the
nonlinear equations on a tail of the sequence. The resulting invariant
island has positive area in each return chart, and hence positive mass
for that billiard's collision probability.

The construction uses centrally symmetric normal graphs over
\(\gamma_E\). Sections~\ref{sec:reflection}--\ref{sec:model} construct
the returns. Sections~\ref{sec:actions}--\ref{sec:restoration} restore the
links, and Section~\ref{sec:entropy} transfers entropy to the collision map.


\section{Reflection at a perturbed wall}\label{sec:reflection}

For a smooth strictly convex table \(D\), use the annulus of oriented lines crossing
its interior. Write \(q_+(\ell)\) for the forward exit of an oriented
line \(\ell\), and \(T_D\ell\) for the new line after reflection there.
The outgoing collision state immediately after reflection at
\(q_+(\ell)\) identifies \(T_D\) with the collision map \(B_D\).
For a centrally symmetric table, the antipodal transformation commutes
with reflection. We write \(\overline T_D\) for the quotient by this
transformation, which acts freely on oriented lines.

\subsection{Fixed charts and parameters}\label{fixed-charts-and-parameters}

Let \(\gamma\) be an arclength arc of the fixed ellipse, \(t=\gamma'\), and
\(n\) its outward unit normal. With positive boundary orientation,
\(t'=-\kappa n,\ n'=\kappa t\). Fix one reference incoming ray whose last
intersection in its direction is \(\gamma(s_*)\), with outward normal component
\(w>0\). Let \(z=(x,y)\) be a fixed smooth canonical chart of the unperturbed
outgoing rays near the reference reflection, with symplectic form
\(dy\wedge dx\). The chart may include the fixed antipodal transformation on the
way from reflection to the representative of the quotient. In the formulas below
that transformation is included in the outgoing chart. The center
\(z_*=(x_*,0)\) is a point on the local stable axis and the matching old impact
function is \(s_b(z)\). The input ray direction \(v(z)\) is obtained by
inverting the fixed reflection (and undoing the antipodal transformation if
needed).

Use independent small parameters \(\mu,\epsilon\) and an additional parameter
\(d\), with \(|d|\leq 1\). On fixed bounded \((X,Z)\)-sets put \[
 z=z_*+\mu(\alpha X,dZ),\qquad f(s)=\epsilon\mu\,P((s-s_*)/\mu)
\] with fixed nonzero \(\alpha\). The physical scaling will set
\(\epsilon=\tau=3^{-k}\), \(\mu=1/k\), \(d=\tau/\mu\), so the height of the
input tube in physical coordinates is \(O(\tau)\). Let \(P\) be smooth on a
common fixed interval with a collar around all sample points, with bounded
\(C^2\) norm. A zero extension is used for supported profiles. The base impact
coordinate \(\sigma_b=(s_b(z)-s_*)/\mu\) is uniformly smooth in the
parameterized input, including \(\mu=0\), by the integral formula for the
difference of \(s_b\) at \(z\) and \(z_*\).

\subsection{The exact impact and readout}\label{the-exact-impact-and-readout}

Write \(s=s_*+\mu\sigma\), \(s_b=s_*+\mu\sigma_b\) and \(\gamma_f=\gamma+fn\).
The ray crosses \(\gamma(s_b)\). The selected perturbed impact satisfies \[
 D(\sigma,\sigma_b,z)(\sigma-\sigma_b)
   +\epsilon\,\det(n(s),v(z))P(\sigma)=0,\qquad
 D=\int_0^1\det(\gamma'(s_b+\theta(s-s_b)),v(z))\,d\theta.
\] The factor \(D\) has fixed nonzero sign and a positive lower bound in
magnitude on the chosen transverse neighborhood. The derivative of the left side
with respect to \(\sigma\) has that same property for small \(\epsilon,\mu\), by
transversality and the reference bound on \(P'\). A fixed compact collar and the
intermediate value theorem give a unique root there with
\(|\sigma-\sigma_b|\le C|\epsilon|\|P\|_0\). This is the actual exit from the
strictly convex table: the incoming direction has positive outward normal
component there, and a line can have only one transverse exit. Define \[
 B=-\det(n(s),v(z))/D,\qquad s-s_b=\epsilon\mu\, B P(\sigma).
\]

The impact point differs from \(\gamma(s_b)\) by \(\epsilon\mu\, U\), where \[
 U=P(\sigma)\left(n(s)+B\int_0^1\gamma'(s_b+\theta(s-s_b))\,d\theta\right).
\] The unnormalized tangent to the perturbed arc is \[
 t_f^{\mathrm{un}}=(1+\epsilon\mu\kappa(s)P(\sigma))t(s)+\epsilon P'(\sigma)n(s).
\] Subtract \(t(s_b)\) using \(s-s_b=\epsilon\mu BP(\sigma)\), producing
\(\epsilon\) times a smooth expression in the fixed coefficients and the values
\(P(\sigma), P'(\sigma)\). The segment from \(t(s_b)\) to \(t_f^{\mathrm{un}}\)
lies in a fixed annulus about the unit circle when \(\epsilon\) is small. Unit
normalization, rotation to the normal, and reflection
\(v\mapsto v-2(v\cdot n_f)n_f\) preserve the exact factor \(\epsilon\) in the
directional increment, by the integral formula for a smooth difference on this
neighborhood.

An oriented line near the reference has smooth coordinates consisting of a local
angle of its unit vector \(u\) and a signed offset \(\det(a,u)\), where \(a\) is
any point on that line. At the output use \(a=\gamma_f(s)\). The position and
direction differences just found imply that the output line-coordinate
difference has the factor \(\epsilon\). Apply the same smooth difference formula
to the fixed outgoing chart from line coordinates to \(z\). Thus the actual
relative reflection \(g_f\), which sends the reference reflected ray to the new
reflected ray from the same incoming line, has the exact identity \[
 g_f(z)-z=\epsilon\,L(\mu,d,\epsilon,X,Z;P).
\] The components of \(L\) are smooth coefficient expressions evaluated at the
root and at \(P(\sigma),P'(\sigma)\). All coefficients extend across zero scale
parameters and have bounded derivatives of each fixed order on compact subsets. The only denominators above come from uniform transversality, the tangent norm,
or fixed coordinate changes. In particular, the differences of the two outgoing
coordinates are taken before dividing by the output affine scales.

\subsection{Finite-order functional
estimates}\label{finite-order-functional-estimates}

Fix \(r,q\geq0\) and a bound on the \(C^{r+q+4}\) norm of the profiles, with
domains and collars fixed. Implicit differentiation of the exact impact equation
in \((X,Z,\epsilon)\) and \(q\) profile directions isolates each highest root
derivative with coefficient \(\partial_\sigma F\), where \(F\) is the left side
of the impact equation. Remaining terms have earlier root derivatives and
evaluations of finitely many jets of \(P\) and the directions at the same root.
Induction through total order at most \(r+q+1\), followed by differentiation of
\(L\), uses at most \(r+q+2\) jets of profiles and directions. This gives
uniform mixed finite-order bounds, including one \(\epsilon\) derivative, by a
constant times the product of the direction norms in \(C^{r+q+4}\). For \(q=0\)
the product is 1. In particular \(L(\epsilon)-L(0)=O_{r,q}(\epsilon)\) in these
bounds. The estimates concern output from the enlarged physical tube with
bounded \((X,Z)\).

When \(\epsilon=\tau\), dividing the nonlinear remainder \(\tau[L(\tau)-L(0)]\)
by the output scales \(\alpha\mu,\beta\tau\mu\) costs at most \(C\tau/\mu\),
with fixed \(\beta\ne0\). The same assertion holds after differentiating with
respect to finitely many bounded profile directions, and after substitution of a
uniformly smooth bounded incoming \((X,Z)\)-parameterization.

\subsection{Sign and comparison with the prescribed
shear}\label{sign-and-comparison-with-the-prescribed-shear}

Use coordinates \((s,p)\) of an outgoing ray at its first intersection with the
unperturbed curve, with the inward-directed branch \(u=p\,t-w n\) and
\(w=\sqrt{1-p^2}\). The incoming direction is \(v=p\,t+w n\). These coordinates
have form \(dp\wedge ds\): the line invariant two-form is the exterior
derivative of \(u\cdot da\), where changing the anchor \(a\) along \(u\) adds an
exact differential. This is the first-intersection counterpart of the
last-intersection chart in \cite[\S1.4]{Glutsyuk}. The antipodal readout
preserves this form.

At amplitude zero the variation of the new impact parameter is
\(\delta s=pf/w\). Indeed the first-order position difference \(\delta s\,t+fn\)
must be parallel to \(v\). The unit tangent varies by \((f'-\kappa\delta s)n\),
and the reflected direction by \(2(f'-\kappa\delta s)(w t+p n)\). The first
intersection of the new reflected line with the old arc has variation
\(\delta\widehat s=2pf/w\), by sliding \((f/w)u\) along the outgoing line to
first order. Consequently its tangential momentum varies by
\(2w(f'-\kappa\delta s)+\kappa w\delta\widehat s=2w f'\). Thus the first
variation of the post-relative reflector is \[
 (-K_p,K_s),\qquad K(s,p)=2w(p)f(s).
\] The sign here uses \(\iota_{X_K}(dp\wedge ds)=dK\). The two curvature terms
cancel.

In the fixed output chart this reads \(K(x,y)=A(x,y)f(s_b(x,y))\), \(A=2w\),
with \(X_K=(-K_y,K_x)\). On the selected stable segment require
\(\partial_x s_b(x,0)\ne0\), and choose \(f(s_b(x,0))=G(x)/A(x,0)\). The shape
definition includes a cutoff in a common scaled \(x\)-interval that is one on
the required larger impact interval. The exact formula \[
 K(x,y)=R(x,y)G(\xi(x,y)),\qquad
 s_b(\xi(x,y),0)=s_b(x,y),\qquad
 R(x,y)=A(x,y)/A(\xi(x,y),0)
\] has \(\xi-x=y E\) and \(R-1=y E_1\), with \(E,E_1\) smooth and all their
fixed-order derivatives bounded. Let
\(G(x_*+\alpha\mu X)=\tau\mu\,\mathcal G(X)\) with uniformly bounded profiles
\(\mathcal G\) in the prescribed finite norms. The corresponding \(P\) above has
the same finite-order uniformity by the inverse footpoint map, since the
rescaled map and its derivative bounds follow from a fixed smooth local inverse
and the identity for a divided increment. In raw coordinates \((X,Z)\) the
argument of \(\mathcal G\) at \(\xi\) differs from \(X\) by
\(O_{C^r}(\tau/\mu)\). Differentiating the last exact formula gives \[
 K_y=O_{C^r}(\tau),\qquad
 K_x-G'(x)=O_{C^r}(\tau^2/\mu).
\] For the second estimate, \(R\xi_x-1=O_{C^r}(\tau)\), \(R_x=O_{C^r}(\tau)\),
and the integral difference of \(G'(\xi)\) and \(G'(x)\) costs
\(C_r\tau(\tau/\mu)\). The estimates use at most \(r+2\) profile derivatives.
The calculation is linear in \(G\), including the cutoff, so it applies to all
first profile directions and vanishes at higher profile orders. Put \(z=(x_*+\alpha\mu X,\tau Z)\), and let
\[
 {\cal A}_{\tau,\mu}(x,y)
   =\left(\frac{x-x_*}{\alpha\mu},\frac{y}{\beta\tau\mu}\right).
\]
The raw pair \((X,Z)\) ranges over a bounded set. Comparing the two
affine outputs at the same raw input gives
\[
 {\cal A}_{\tau,\mu}(g_f(z))-{\cal A}_{\tau,\mu}(z)
  =\left(0,\frac{G'(x)}{\beta\tau\mu}\right)+O_{C^r}(\tau/\mu^2).
\]
The term on the right is the leading vertical kick. The horizontal
component of the linearized difference is \(O_{C^r}(\tau/\mu)\).
The error bound holds through any fixed number of profile derivatives
on bounded sufficiently high norms, with the corresponding multilinear
direction norms.

\subsection{Uniform tame control of free
increments}\label{uniform-tame-control-of-free-increments}

A scalar factorization provides a stronger local fact about the unknown free
profiles on a fixed low-norm neighborhood. Write \(P=P_k^0+\mu U\) with
\(P_k^0\) a fixed uniformly smooth base family and \(\|U\|_{C^2}\)
bounded. The scaled fixed footpoint reparametrization and smooth cutoff are
uniformly bounded linear operations on each \(C^r\), so \(U\) may be the
rescaled free profile from the physical formula above. Directly from the exact
reflection identity, \[
 L(P_k^0+\mu U)-L(P_k^0)
   =\mu\int_0^1 D_P L(P_k^0+t\mu U)[U]\,dt .
\] After normalization the vertical difference is the second component of this
integral divided by \(\beta\). The horizontal difference is the first component
times \(\tau/\alpha\). This treats the free correction at a common raw input and
keeps the cancellation before division by \(\tau\mu\).

In an expression formed by differentiating the root and readout formulas, assign
weight \(\max(j-1,0)\) to a jet \(P^{(j)}(\sigma)\), and \(j\) to a direction
jet \(U_i^{(j)}(\sigma)\). All smooth fixed coefficients, denominators and the
unvaried \(P,P'\) are weight zero and uniformly bounded. Keep denominators as
powers of the scalar impact derivative, the tangent norm or a smooth
fixed-coordinate factor. The first state derivative of the root is a weight-zero
expression, and its derivative in one new direction is a weight-zero coefficient
times \(U_i(\sigma)\). These two facts come from dividing \(-F_{a}\) or
\(-D_PF[U_i]\) by \(F_\sigma\) in the impact equation. A subsequent state
derivative can increase the total weight by at most one. A derivative in a new
profile direction can substitute \(U_i^{(j)}\) for \(P^{(j)}\), or add the root
derivative in a composition, or differentiate a denominator. Each increases the
weight by at most one and introduces the direction exactly once. This remains
true for one amplitude-parameter derivative, which has the same root bound as an
ordinary state derivative. By induction every term of \(D_{X,Z}^rD_P^q L\) has
total jet weight at most \(r+q\).

For \(m=r+q>0\), the finite-interval interpolation inequality bounds a \(P\)-jet
of weight \(j-1\) using low \(C^1\) norm and high \(C^{m+1}\) norm, with high
exponent \((j-1)/m\). For a direction jet it uses low \(C^0\), high \(C^m\) and
high exponent \(j/m\). This inequality follows by taking an appropriate \(j\)-th
Taylor difference on a one-sided interval of length comparable to \(h\) and
balancing \(h^{-j}\|u\|_0\) against \(h^{m-j}\|u\|_{C^m}\), applied to \(u=P'\)
or \(u=U_i\). The sums of high exponents are at most one. The weighted
arithmetic-geometric mean inequality then bounds each term by a constant times
\[
 (1+\|P\|_{C^{m+1}})\prod_i\|U_i\|_{C^0}
 +\sum_i\|U_i\|_{C^m}\prod_{j\ne i}\|U_j\|_{C^0}.
\] A zero low norm makes its directional term identically zero. The case \(m=0\)
follows from the undifferentiated coefficient bound. The bounds are uniform on a
fixed low-norm profile neighborhood, for each \(r,q\), with loss independent of
\(r\).

The same argument applies under the integral for the free increment. Its
\(q\)-th \(U\) derivative for \(q\geq1\) has a prefactor \(\mu^{q-1}\) in the
vertical coordinate from the difference quotient. For \(q=0\) use the integral
itself, increasing the input shift by at most one. If the base family has
uniformly bounded higher norms, these estimates give local smooth tame bounds
for the free increment with a fixed shift for every fixed \(q\).

\subsection{Normalized reflection
estimate}\label{normalized-reflection-estimate}

On a compact transverse old reflection chart with a regular stable-footpoint
coordinate, an \(O(\tau\mu)\) supported physical primitive determines a true
locally perturbed exit and reflection. Its normalized output on
\(|x-x_*|\leq C\mu,\ |y|\leq C\tau\) equals the prescribed vertical kick up to
\(O(\tau/\mu^2)\) in each fixed finite state/profile norm with finite input
loss. For \(G=G^0+\alpha\beta\tau\mu^2 V\) the free increment has the uniform
fixed-loss local estimates above. Here \(\tau=3^{-k}\), \(\mu=1/k\),
\(\alpha,\beta\ne0\) are the two fixed affine factors, and \(V\) is the
fixed-domain smooth primitive.

\section{Realization of a finite Hénon word}\label{sec:word}

\subsection{Finite-word realization}\label{finite-word-realization}

Fix an odd integer \(N\ge1\), a smooth primitive profile \(V_i\) for each factor
\(H_{V_i'}(X,Y)=(Y,-X+V_i'(Y))\), and a compact starting set \(K\). Fix compact
partial domains with positive margins containing the partial images of \(K\) and
common bounds on the relevant profile intervals. These choices precede the limit
parameter. The claim is that for \(k\) sufficiently large there is one smooth
strictly convex centrally symmetric table given by an explicit normal graph on
\(E=(2,\sqrt3)\), and fixed-for-\(k\) affine charts in the old oriented-ray
coordinates, such that:

\begin{enumerate}
\def\labelenumi{\arabic{enumi}.}
\tightlist
\item
  Each starting ray follows \(N\) consecutive prescribed blocks, each with \(k\)
  unchanged saddle iterates followed by a fixed finite transit. Only the last
  impact of each transit can use a changed part of the table, at its own
  designated window.
\item
  The normalized \(N\)-block return converges to the given word with joint state
  and profile derivatives of every prescribed finite order. More precisely, for
  any fixed \(r,q\), convergence in \(C^r\) of the state derivatives of
  \(D_V^j\), \(0\le j\le q\), is uniform when profiles and their \(C^{r+q+6}\)
  norms are bounded, the strict margins are uniform, and the direction norms are
  at most 1 in \(C^{r+q+6}\). We work on a fixed small neighborhood of the input
  compact where the same margins hold.
\item
  The boundary displacement is \(O_{C^r}(3^{-k}k^{r-1})\) for fixed \(r\) and a
  bounded \(C^r\) family of primitives. The affine charts have a common, nonzero
  signed Jacobian proportional to \(3^{-k}k^{-2}\).
\end{enumerate}

Here the profile parameters may be taken on the fixed compact intervals extended
with smooth collars.

\subsection{Physical coordinates and
support}\label{physical-coordinates-and-support}

Use the quotient by central symmetry introduced in
Section~\ref{sec:reflection}. Its area form descends from \(dp\wedge ds\).
In the ellipse write \(\ell\) in collision coordinates \((\theta,v)\)
at its current forward exit \(q_+(\ell)=\gamma_E(\theta)\). Here
\(p\) is the tangential component at that impact and
\(v=|\gamma_E'(\theta)|p\). With \(s_E\) the ellipse arclength,
\(dv\wedge d\theta=dp\wedge ds_E\).
The involution sends \((\theta,v)\) to \((\theta+\pi,v)\).

The unperturbed focal separatrix has parametrization \[
 q(t)=(2\tanh t,\sqrt3\,\mathrm{sech}\,t).
\] Here \(q(t)\) is a boundary exit on the upper half of the ellipse,
and the physical successor is \(-q(t+\log3)\). The quotient advance
is \(t\mapsto t+\log3\). The expanding and contracting
multipliers are 3 and \(1/3\) on the quotient. The separatrix formula follows
from \cite[Lemma 4.1]{DRR}, and the focal chord identification follows from
\cite[Lemma 2.2]{DCRR}. In canonical angle-momentum coordinates \((\theta,v)\),
the two tangent lines are \(v=\pm\theta\). The stable sign at the selected
major-axis point is \(+\), since \(q(t)\) tends to that point as \(t\to+\infty\)
with \(v=\sin\theta\). Fix canonical saddle coordinates \((x,y)\) with stable
axis \(y=0\) and unstable axis \(x=0\), and differential
\(\theta=(x-y)/\sqrt2,\ v=(x+y)/\sqrt2\). Choose the axes to have linear
restrictions with multiplier \(\lambda=1/3\) and \(1/\lambda\), as in
\cite[\S6.1]{BT}. Canonical straightening follows by straightening the
transverse axes, then integrating the density of the area form in the second
variable. Linearization on each one-dimensional axis follows from the limit
\(\lambda^{-j} f^j\) for a smooth contraction \(f\) with derivative \(\lambda\)
at zero, using \(f^{-1}\) on the unstable axis. The limit is smooth on a small
interval because telescoping increments are \(O_{C^m}((\rho^2/\lambda)^j)\) with
\(\lambda<\rho<\sqrt\lambda\). Here the bounds \(\|f^j\|_{C^m}=O(\rho^j)\)
follow inductively by differentiating the recurrence, with the term containing
the highest derivative multiplied by \(\sup|f'|<\rho\) and all products of lower
positive-order derivatives \(O(\rho^{2j})\). The canonical extensions
\((x,y)\mapsto(h(x),y/h'(x))\) and \((x,y)\mapsto(x/j'(y),j(y))\) preserve the
other axis restriction when \(h'(0)=j'(0)=1\).

The current exit in this chart has arclength
\(s_E(\theta(x,y))\) up to the fixed antipodal representative.
For the relative reflection of Section~\ref{sec:reflection},
\(s_b(x,y)\) is the forward exit of the inverse old ray, evaluated
at an old output \((x,y)\). On the stable axis,
\[
 s_b(x,0)=s_E(\theta(3x,0))
\]
in the corresponding arclength chart. Its derivative at the vertex
has magnitude \(3\sqrt3/\sqrt2\). This supplies the regular arclength
coordinate for the control profiles. The saddle box is protected
using its \emph{current} exits.

Choose a small positive stable input center \(x_c\) for a control impact. A
saddle box \(U\) of radius \(u\in(0.37x_c,0.43x_c)\) contains the old image of
that stable input, at \(x_c/3\), and excludes the input. The whole old
last-footpoint projection of \(U\) is at arclength distance at most
\(2u\sqrt3/\sqrt2+O(u^2)\) from the vertex. The projection of a small stable
center interval around \(x_c\) is at least \(0.98x_c\sqrt3/\sqrt2+O(x_c^2)\)
away. This follows from the differential of the chart and
\(|\gamma_E'(0)|=\sqrt3\). The stated radii give strict separation for
sufficiently small fixed \(x_c\). Shrink the center interval so it lies in one
compact separatrix time interval of length less than \(\log3\). Choose the box
radius to avoid an exact unstable boundary hit along the chosen reference orbit.
The excluded radii along this one discrete orbit are countable. After shrinking
the center interval again, one common finite \(m\) takes its last unstable exit
points inside \(U\) to the stable landing points inside \(U\), with all
intermediate points outside \(U\). Indeed near each end of the homoclinic orbit
one coordinate on the appropriate axis changes by \(\lambda\) or \(1/\lambda\)
until leaving the small chart. The remaining compact portion of the separatrix
is disjoint from the saddle and stays outside sufficiently small boxes. All
inequalities retain strict margins on a smaller center interval.

Take \(N\) distinct control centers there. Let \(M_i^+=(x_i^+,0)\) be their old
outputs, and let \(M_i^-=(0,y_i^-)\) be the unstable starting point of the
associated \(m\)-step transit. Apart from the designated last input, the finite
native transit impacts lie outside all control arcs with strict physical
arclength separation. The native impacts have times shifted by nonzero integer
multiples of \(\log3\) relative to the selected impact. Their physical
footpoints modulo the antipodal copies are separated by injectivity of \(q\) on
the open upper semi-ellipse and by the common time window of length less than
\(\log3\). Thus the same separation holds on fixed neighborhoods of each finite
transit up to the last reflection. One can additionally separate the designated
input arcs from each other. For \(k\) large the \(O(\mu)\) supports lie inside
these fixed arcs. If an old exit point is off the supports, the point and its
normal are unchanged. For a convex perturbation with the same crossing at that
location it is still the unique forward exit on the line, so that reflection is
exactly the old one.

\subsection{Unperturbed long saddle passage inside the same
box}\label{unperturbed-long-saddle-passage-inside-the-same-box}

For the fixed old map in the saddle chart write
\(\bar x=\lambda x+p(x,y)x,\ \bar y=\lambda^{-1}y+q(x,y)y\). The axis
linearizations give \(p(x,0)=q(0,y)=0\). The area identity on each axis gives
\(p(0,y)=q(x,0)=0\) by integrating its first derivative there. Thus the error
coefficients \(p,q\) both have a smooth factor \(xy\). On a small box the mixed
equations for a step can therefore be written as \[
 x_{j+1}=\lambda x_j(1+x_j y_{j+1}a(x_j,y_{j+1})),
 \qquad y_j=\lambda y_{j+1}(1+x_j y_{j+1}b(x_j,y_{j+1})).
\] In the second equation invert the old \(\bar y\) for a fixed \(x_j\), using
its uniformly positive derivative near \(\lambda^{-1}\). Both correction
coefficients are smooth and uniformly bounded with fixed-order derivatives on a
small enlarged box.

For \(k\) steps set \(A_j=\lambda^{-j}x_j,\ B_j=\lambda^{-(k-j)}y_j\) with
endpoints \(A_0=x_0,\ B_k=y_k\). The factor \(x_jy_{j+1}\) is
\(\lambda^{k-1}A_jB_{j+1}\). Sum the recurrences for \(A_{j+1}-A_j\) and
\(B_j-B_{j+1}\) to express the entire intermediate arrays as a fixed point. The
right side equals the constant arrays \(x_0,y_k\) plus sums of at most \(k\)
terms with factor \(\lambda^{k-1}\). On the sup-norm box of radius \(2u\) in the
arrays it has Lipschitz constant \(O(k\lambda^{k-1})\). Each derivative of any
prescribed finite order is bounded by \(C_r k\lambda^{k-1}\) for the correction,
since rescaling the arguments costs only factors of magnitude at most 1. For
large \(k\) this defines the unique solution in that box, with higher endpoint
derivatives of the same small order for the correction, by differentiation and
the uniform inverse of the linearized fixed point equation. In the physical
coordinates every intermediate component lies in \(U\), since the endpoint
magnitudes are less than \(u\) and \(2u\lambda<u\) for the other entries. In
fact \(|x_j|\le \rho^j|x_0|\) with fixed \(\rho<0.5\) on the small saddle box.
This argument establishes the old local passage on the actual box separated from
the physical perturbations. The physical perturbation preserves each of these
steps exactly.

In particular, \[
 x_k=\lambda^k(x_0+\xi_k(x_0,y_k)),\qquad
 y_0=\lambda^k(y_k+\eta_k(x_0,y_k)),
 \qquad
 \|(\xi_k,\eta_k)\|_{C^r}\le C_r k\lambda^{k-1}.
\] The proof uses only one fixed unperturbed saddle map.

\subsection{Rescaling, synchronization and
induction}\label{rescaling-synchronization-and-induction}

Write the fixed outer transit from \(M_i^-\) to \(M_i^+\) near its endpoints as
\[
 x'=x_i^++b_i(y-y_i^-)+R_i^{(1)}(x,y-y_i^-),\qquad
 y'=x\big(c_i+d_i(y-y_i^-)+R_i^{(2)}(x,y-y_i^-)\big).
\] Both remainders are \(O(|x|+|y-y_i^-|^2)\) and are smooth, as follows by
Taylor expansion and the fact that the unstable axis maps to the stable axis on
this interval. At the center the determinant is \(-b_ic_i=1\), so \(b_i,c_i\)
are nonzero. All \(b_i\) have a common sign by continuity on a sufficiently
small interval.

Use cyclic indices and choose fixed positive \(\rho_i\) by
\(\rho_i=-c_i b_{i-1}\rho_{i-2}\). Since \(N\) is odd, the step of two goes
through the entire cycle, and the product of all ratios
\(-c_i b_{i-1}=b_{i-1}/b_i\) is one. Thus a positive choice exists with one
scale fixed freely, including \(N=1\). Set \(a_i=b_i\rho_{i-1}\),
\(\tau=\lambda^k\), \(\mu=1/k\), and \[
 Q_{ik}(X,Y)=(x_i^++a_i\mu X,\ \tau(y^-_{i+1}+\rho_i\mu Y)).
\] The signed Jacobian \(a_i\rho_i\tau\mu^2\) is independent of \(i\) by the
recursion. We use zero extra centering corrections since the mixed saddle
estimate above gives \(k\lambda^{k-1}/\mu\to0\).

For a bounded normalized input at slot \(i-1\), the long saddle passage ends at
\[
 \tilde x=\tau(x_{i-1}^++a_{i-1}\mu\widetilde X),\qquad
 \tilde y=y_i^-+\rho_{i-1}\mu\widetilde Y,
 \qquad (\widetilde X,\widetilde Y)=(X,Y)+o_{C^\infty}(1).
\] Indeed the exact mixed equation for \(\tilde y\), after division by \(\mu\),
has an error \(O_{C^r}(k\lambda^{k-1}/\mu)\) for every \(r\) on bounded
normalized sets and can be solved with derivative uniformly close to the
identity. The formula for \(\tilde x\) has the same error. The fixed reference
transit then gives \[
 X_{\rm raw}=\widetilde Y+e_{1k},\qquad
 Y_{\rm raw}=C_i/\mu+A_i\widetilde Y-\widetilde X+e_{2k},
\] where \(C_i=(c_i x_{i-1}^+-y_{i+1}^-)/\rho_i\) and
\(A_i=d_i x_{i-1}^+\rho_{i-1}/\rho_i\). The \(-\widetilde X\) coefficient is
\(c_i a_{i-1}/\rho_i=-1\). Both errors tend to zero in each \(C^r\) norm by the
stated smooth Taylor remainders after substitution and the fact that
\(\tau/\mu\to0\). In particular the physical raw output can be written as
\(x=x_i^++a_i\mu X_{\rm raw}\), \(y=\tau Z_{\rm raw}\), where
\((X_{\rm raw},Z_{\rm raw})\) ranges over a bounded family with all fixed-order
derivatives bounded. The potentially unbounded coordinate \(Y_{\rm raw}\) is
never used as a bounded input of a local reflection lemma.

At slot \(i\), prescribe the derivative of the primitive using the raw
horizontal coordinate \[
 G_{ik}'(x_i^++a_i\mu X)=\tau\rho_i[-C_i+\mu(-A_i X+V_i'(X))]
\] on a scaled interval with a strict collar around the relevant
\(X_{\rm raw}\). Integrate and use a smooth cutoff that is one on that interval
and vanishes beyond a slightly larger fixed scaled interval. In this coordinate
its primitive is \(\tau\mu\) times a uniformly smooth fixed base term plus
the free term \(a_i\rho_i\tau\mu^2 V_i\) with the cutoff (an additive constant
of the free primitive can be fixed). With the \(s_b\) and \(2w\) of the
reflection calculation, use \(f_{ik}(s_b(x,0))=G_{ik}(x)/(2w(x,0))\). The fixed
map \(s_b(x,0)\) here is the last-footpoint map on the inverse old stable
branch, so its derivative is nonzero. All \(N\) supports can be made disjoint
for large \(k\), including their antipodal copies. Reflecting at one does not
evaluate any other profile. Their sum, extended with central symmetry, gives the
single table. For each control slot the linear variation of this very wall is
\(J_k(\chi_i U_i)(\sigma(s))/(2w_{\rm old}(s))\) on its support, with the
analogous antipodal copy. Here \(\sigma(s)\) is its fixed scaled inverse old
stable-footpoint coordinate, \(\chi_i\) is the fixed primitive cutoff, and
\(J_k=a_i\rho_i\tau\mu^2\) is that row's signed Jacobian. The other
profile-dependent terms live on disjoint physical supports. Local inverse and
product estimates in scaled arclength give the claimed \(C^r\) boundary bounds.
In particular the boundary stays embedded and strictly positively curved, by
\(C^2\) smallness and the strict curvature of the ellipse.

Each finite transit is exact until its designated last reflection, by the fixed
physical margins. For the latter use the true exit and the estimate in
Section~\ref{sec:reflection}, with \(\beta=\rho_i\). After
normalization the kick is \(-C_i/\mu-A_iX_{\rm raw}+V_i'(X_{\rm raw})\) plus an
error tending to zero in all the fixed finite mixed norms. Adding it to
\(Y_{\rm raw}\) cancels the constant exactly, and produces \[
 (X_{\rm raw},-\widetilde X+V_i'(X_{\rm raw})+e_{2k}-A_i e_{1k})
\] up to that small error. This converges jointly to \(H_{V_i'}(X,Y)\). Both the
intermediate reference estimates and the final profile dependence are evaluated
on common bounded physical or normalized domains as specified in their
derivation.

Induct over the fixed \(N\) slots. A normalized output close to the
corresponding ideal partial image lies in the next compact domain with its
margin. Its physical image is in \(U\) near that slot's \(M_i^+\), hence the
next saddle block is the allowed one. This proves the itinerary for every ray in
the input compact, uniformly with the given margins and profile norms. For mixed
derivatives up to total order \(m=r+q\), apply the finite chain and product
rules in the state and the profile directions. Each block uses at most \(m\)
total derivatives. Comparing evaluations at actual and ideal partial images uses
one extra state derivative of the limiting factor on the common enlarged domain.
The finite composition converges in these mixed norms on a bounded \(C^{m+6}\)
family. The distinct landing centers lie in a narrow stable interval near
\(x_c/3\). For \(k\) large every intermediate local step contracts \(|x|\) below
that entire interval, and intermediate outer steps lie outside \(U\). Thus the
initial landing chart is visited again only after the cyclic \(N\)-block return,
even if a smaller starting compact is used.

\section{The reference island and diagonal approximation}\label{sec:model}

\subsection{The marked reference island}\label{the-marked-reference-island}

We use the standard-area map \(\check F\) from
\cite[Proposition 2.1 and \S3.4]{BT}. Its comparison domain is the closed disc
\(D_{\mathrm{BT}}\), and it has a smooth plane extension
\(\check F_{\mathrm{ext}}\) equal to the identity outside that disc. A bi-link
here consists of two invariant heteroclinic arcs from the same unstable endpoint
\(P\) to the same stable endpoint \(Q\), forming a smooth embedded circle with
the endpoints. The reference map has one outer circle and three disjoint inner
circles in its interior. Their eight endpoints are distinct hyperbolic fixed
points. The region inside the outer circle and outside the inner closed discs is
a stochastic island in the sense of \cite{BT}.

Let \(\Phi\) denote the area-preserving preparation constructed below
and \(F_*=\Phi\check F_{\mathrm{ext}}\Phi^{-1}\) the model map.
On each prepared circle, write \([x_a-2h,x_a]\) and
\([x_b,x_b+2h]\) for the horizontal projections of the two strips,
where \(h>0\) is fixed for that circle.

Use smooth rounded rectangles with horizontal upper and lower straight pieces as
the four model bi-link circles. Give the three inner rectangles disjoint
horizontal projection intervals inside the outer rectangle, leaving a fixed
interval to their left. The possibility follows from the strict area inequality:
the outer disc area exceeds the sum of the three inner disc areas. For example,
before rounding, choose a common height \(H\) for inner rectangles less than the
outer height \(H_0\), close enough to \(H_0\) that the sum of their widths
\(A_i/H\) is strictly less than \(A_0/H_0\). Use this remaining width for the
left interval and positive gaps. Rounding close to the corners and adjusting the
widths slightly preserves the exact prescribed areas, the gaps, and the straight
pieces. Use the left interval for the two outer strips and both outer saddle
projections. For each rectangle choose its \(a\)-strip to the left of its
\(b\)-strip. On the upper straight segment put \(P\) just to the right of the
\(a\)-strip and \(Q\) just to its left, both strictly outside the two strips.
Then \(L_a\) is the short upper segment directed from right to left, and \(L_b\)
is the other arc, first crossing the \(b\)-strip rightward on the top and later
crossing it leftward along the bottom. The widths and strip locations can be
chosen with strict margins, independently for the four rectangles. This gives
the labelled circles and straight portions for the local construction of
\cite[Lemma 4.5]{BT}. For the first leading fundamental interval and its
translate, the upper coordinate is the horizontal phase \(t\), with area form
 \(de\wedge dt\). On the upper \(b\)-strip these coordinates are
\(t=x_b+hu\) and \(e=\epsilon/h\) in the time-energy chart used below.
The prepared map translates by \(-h\) along the \(a\)-branch
and by \(+h\) along the \(b\)-branch, including collars of the first interval
into the second. For the \(b\)-branch use also the time-energy chart
\((u,\epsilon)\) with unit time step. The coordinate maps in the proof of
\cite[Lemma 4.5]{BT} are \[
 (u,\epsilon)\longmapsto(x_b+hu,y_b+\epsilon/h),\qquad
 (u,\epsilon)\longmapsto(x_b+4h-hu/2,y_b'-2\epsilon/h).
\] The first formula applies on the upper two-period interval \(0\le u\le2\) and
the second on the lower four-period interval \(4\le u\le8\), with the retained
collars. Here \(h>0\) is the horizontal step of the upper chart, \(x_b\) is the
left edge of the \(b\)-strip, and \(y_b,y_b'\) are its upper and lower heights.
The fourth iterate of a point starting at horizontal phase \(t\) on the first
upper interval has horizontal coordinate \(x_b+2h-(t-x_b)/2\). Each of the next
three steps subtracts \(h/2\). All eight endpoint projections avoid all eight
closed strips. Each strip also misses each of the other inner circles by a
positive margin.

For the extension step, take local symplectic annular maps from the construction
in \cite[Lemma 4.5]{BT}. The source and target annuli surround discs of pairwise
matching areas. Extend the orientation-preserving maps on smaller collars to a
smooth plane diffeomorphism \(\Phi_0\), identity outside a large ball. One can
extend component by component using the smooth classification of compact planar
surfaces with boundary and the extension of orientation-preserving boundary
diffeomorphisms over a collar. The complementary compact components are discs
and a disc with three holes, and the exterior is completed inside the large ball
with its outer collar fixed. On each complementary component the compactly
supported difference \(\Phi_0^*\omega-\omega\) integrates to zero, by the
enclosed areas and the equality of the local annular areas. To obtain a
compactly supported primitive of the difference, cover its support and
connecting paths by finitely many overlapping interior coordinate rectangles. A
partition of unity and transfer of the integral over a spanning tree using
unit-integral bumps in the overlaps makes each rectangle's contribution have
integral zero. For \(d(x,y)\,dx\wedge dy\) of integral zero in a rectangle,
choose an interior bump \(\chi(x)\) of integral one, put
\(b(y)=\int d(x,y)\,dx\), and use the primitive \[
 \left[\int_{-\infty}^x (d(v,y)-\chi(v)b(y))\,dv\right]dy
 -\chi(x)\left[\int_{-\infty}^y b(v)\,dv\right]dx .
\] It is supported in the rectangle. The sum is a primitive \(\beta\) supported
away from the collars. The positive interpolated area forms
\(\omega_u=(1-u)\omega+u\Phi_0^*\omega\) give the compactly supported vector
field determined by \(\iota_{X_u}\omega_u=-\beta\). Its time-one map pulls
\(\Phi_0^*\omega\) back to \(\omega\) and equals the identity on the collars, so
composing it with \(\Phi_0\) supplies the required symplectic extension. The
branch labels \(a,b\) on each source circle may be chosen in its two possible
directions to agree with the prescribed plane orientation and order in Lemma
4.5.



\subsection{Choosing the words and the
compression}\label{choosing-the-words-and-the-compression}

Choose one large compact model domain with a collar containing a rectangle whose
interior contains the prepared closed source disc. Include the buffered saddle
boxes and the finite reference phase patches and their images used below. This
choice precedes any word approximation. By Turaev's approximation theorem in the
planar form used in the proof of \cite[Corollary 6.3]{BT}, a smooth symplectic
plane map can be approximated on a fixed compact domain in any finite \(C^r\)
norm by an even number of maps \((X,Y)\mapsto(Y,-X+\psi(Y))\). See also
\cite[Theorem 2]{Turaev}. Apply this to \(H_0^{-1}\circ F_*\) and append
\(H_0\). Choose the resulting odd words \(W_j\) with error at most \(2^{-j}\) in
\(C^j\) on the common collared domain. The finite intermediate compacts may
depend on \(j\).

Append the action-calibration insertion of
Section~\ref{sec:actions} and one control block for each of the eight
strips. Each control block has the form \(H_{\widehat V'}R^3\), where
\(R=H_0\) and \(\widehat V\) is the active primitive multiplied by its
fixed strip cutoff. The wall construction uses its own cutoff, equal
to one on the required impact interval. These are a fixed number of
blocks, each with even
total length. The ideal action-calibration insertion is the identity at zero
calibration, and the ideal free-control blocks are the identity at zero free
profiles. Use the finite-word construction with \(W_j\) and those appended
blocks held fixed, then choose \(k=k_j\). On the suffix all variable reflection
slots are fixed in number. The prefix of the actual return, through \(W_j\)'s
slots, is exactly independent of those variables on the chosen input domain, by
separation of the intervening physical impact sites. Its normalized map
converges to \(W_j\) as \(k\to\infty\), so \(k_j\) can enforce an error of at
most \(2^{-j}\) in every required fixed norm through order \(j\). The closed
words remain odd.

One can choose the finite center sets in increasing order on one fixed small
compact stable arc. The suffix centers can be fixed in a subarc to one side of
the prefix. The outer map and its nonzero transverse derivative \(b(t)\) there
are fixed. If the scale recursion is written \(\rho_i/\rho_{i-2}=b_{i-1}/b_i\),
then the oscillation of \(\log\rho_i\) across the entire cycle is bounded by the
total variation of \(\log|b|\) over the ordered center set, including the one
cyclic wrap. Thus the suffix affine factors can be bounded above and away from
zero uniformly over the varying word lengths. The support widths are reduced
after fixing the center set and are always disjoint. The intermediate prefix
images may require its own larger compacts, which affects the choice of \(k_j\)
but adds no parameter derivative to the suffix.

On the fixed suffix compacts the normalized reflection formulas of
Section~\ref{sec:reflection} use the same fixed-order loss for
each profile derivative. Bounded center data, bounded affine factors, and a
fixed number of variable stages give uniform coefficient estimates there. One
may also impose \(C^j\) bounds on the parameter-independent normalized matrices
and reference remainder for the \(j\)-th row when selecting \(k_j\). For each
prescribed fixed order, the finitely many smaller \(j\) only change the
corresponding uniform constant. The initial prefix approximation is independent
of every variable profile, so it preserves this uniformity when composed in.
This is the mechanism that allows \(N_j\) to grow without a derivative loss
growing with \(j\).

The selected \(k_j\) can also make the boundary change bounded by \(2^{-j}\) in
norms through order \(j\) on fixed bounded profile families. For smooth
corrections that vanish in every fixed norm, the exact normal-graph formula then
ensures convergence to the one ellipse. More explicitly, a bound with the
row-dependent shape coefficient and unit profile norm can be enforced separately
in each order through \(j\). Multiplying by the resulting bounded seminorms of
the solved profiles preserves convergence.

The actual base returns converge to \(F_*\) in every fixed finite order. For
bounded calibration parameters the convergence is joint in state and parameters,
toward \(F_*\) followed by the model action-calibration insertion. This follows by the
finite chain rule on the collared domains along the fixed suffix. The same
convergence to \(F_*\) holds after freezing parameters tending to zero in each
fixed seminorm. Each row has its own \(k_j\) and one physical table.

\section{Periodic actions and protected neighborhoods}\label{sec:actions}

\subsection{Exact ray-space primitive}\label{exact-ray-space-primitive}

In local oriented-line coordinates, choose the point \(a(\ell)\) on the line
\(\ell\) nearest the origin and let \(u(\ell)\) be the unit direction. Define
the globally available line one-form \(\alpha=u\cdot da\). If an anchor on the
same line is \(q=a+t u\), then \(u\cdot dq=\alpha+dt\). The nearest-origin
anchor and the line parameter are invariant in the required sense under the
antipodal transformation, so the form descends to the central quotient.

Let \(T\) reflect the incoming line \(\ell\) at its forward exit \(q\), and
write \(q=a(\ell)+t_{\rm in}u(\ell)=a(T\ell)+t_{\rm out}u(T\ell)\). Equality of
tangential momenta at the collision, with \(dq\) tangent to the fixed table,
gives \[
 T^*\alpha-\alpha=d(t_{\rm in}-t_{\rm out}).
\] Summing this primitive for \(n\) steps gives an action primitive for \(T^n\).
At a quotient fixed point of \(T^n\), the value is the physical sum of the \(n\)
chord lengths. To see this, pair each \(t_{\rm out}\) with the \(t_{\rm in}\) of
its next collision on the same oriented line. The length is the latter minus the
former. The quotient action uses these \(n\) steps even when the lift is
antipodally twisted. Traversing \(2n\) physical reflections always closes the
lift and doubles that value, also for a lift which already closes after \(n\).

Let \(Q\) be the landing chart in a given row and \(A_n\) the ray
action for \(n\) reflections, descended to the quotient. On the
common contractible chart domain, \(Q^*\alpha=JY\,dX+d\chi\), where
\(J\ne0\) is the signed affine scale. For
\(F=Q^{-1}\overline T_D^{\,n}Q\), the primitive of
\(F^*(Y\,dX)-Y\,dX\) is
\[
 S(z)=\frac{A_n(Qz)-\chi(Fz)+\chi(z)}{J}.
\]
At fixed points the two values of \(\chi\) cancel, so each normalized
action difference is the corresponding physical length difference
divided by \(J\).

\subsection{Calibration of the action
differences}\label{calibration-of-the-action-differences}

Write \(P_i,Q_i\), \(0\le i\le3\), for the distinct endpoints of the prepared
stochastic bi-links. Pick a real \(t\) such that the first coordinates of
\(U_tP_i,U_tQ_i\), with \(U_t(X,Y)=(X+tY,Y)\), are pairwise distinct. For pairs
with unequal second coordinates at most one \(t\) is excluded, and pairs with
equal second coordinates already differ in first coordinate. Pick fixed small,
pairwise disjoint collars of these eight projections. Let \(b_i\) be a smooth
supported primitive equal to 1 near the projected \(P_i\) and 0 near all other
seven projections, with support contained in its collar. Insert after \(W_j\)
the action-calibration block \[
 U_{-t}\, S_{(\sum_i \zeta_i b_i)'}\, U_t,\qquad S_\psi(X,Y)=(X,Y+\psi(X)).
\] The symbol \(\zeta_i\) here is a calibration parameter. Each \(S_{\phi'}\) is
\(H_{\phi'}R^3\) for \(R=H_0\). The horizontal shear is a conjugate by \(R\) of
a vertical linear shear. This uses an even number of factors, fixed independent
of the approximating word. Zero free-control blocks can follow it.

For the model family the saddles of \(W_j\) lie in the collars for large \(j\),
and the insertion equals the identity near their images for every small \(\zeta\).
The primitive of \(S_{\phi'}\) relative to \(Y\,dX\) is \(\phi(X)\). Exactness
for composition sums the pulled-back primitives. The \(U_t,U_{-t}\) terms cancel
at the fixed points, where the shear is identity, so the difference
\(\Delta_i=\mathrm{action}(P_i)-\mathrm{action}(Q_i)\) changes by exactly
\(\zeta_i\). The action differences for \(F_*\) vanish because integrating its
action differential along an invariant heteroclinic arc equals the difference
between the integrals of \(Y\,dX\) along the arc and its image. These oriented
arcs are identical.

Joint map and parameter convergence of the actual family on a fixed compact
calibration ball and on a fixed connected model neighborhood gives joint
convergence of these primitive differences. Indeed the one-form
\(F^*(Y\,dX)-Y\,dX\) converges with the joint first derivatives needed here, its
normalized primitive is its line integral from a fixed reference (one can use a
fixed larger disc), and the saddle continuations depend smoothly on the
parameters by hyperbolicity. The exact ray-primitive identification makes this
argument insensitive to the shrinking physical action scale. In the diagonal
construction the actual action vector satisfies \(\Delta_j(0)\to0\),
\(\|D_\zeta\Delta_j-I\|\to0\) uniformly on a fixed small ball. The map
\(\zeta\mapsto\zeta-\Delta_j(\zeta)\) is a contraction and maps the closed ball into itself
for all large \(j\). Its fixed point \(\zeta^{(j)}\) tends to zero and gives exact
equality of each pair's actions. Smoothness of the resulting table is preserved.
This conclusion requires the fixed calibration blocks and joint finite-order
convergence, as supplied by the one-table construction.

\subsection{Protected neighborhoods}\label{protected-neighborhoods}

At each stage fix the calibrated parameters and the earlier
vanishing sequences of repaired profiles, with later profiles zero.
The baseline sets the current active profile to zero. In the model,
the four-factor blocks at zero free profiles have model composite
identity and each nontrivial slot observes the same model \(X\) at its own
four-factor block's output. Each active support lies strictly within its strip.
For large \(j\), actual normalized old footprints in the designated slots are
close to the model \(X\): the inverse old footpoint chart of
Section~\ref{sec:reflection} and the normalized block estimates
give this convergence on bounded inputs. On a fixed small neighborhood of each
saddle this leaves a strict physical support gap after rescaling. Applying the baseline map until the first possible changed reflection and inducting over the
suffix slots proves exact equality of all these nearby return itineraries under
free controls. The saddles, their local maps and their full quotient actions are
therefore unchanged during these variations. The same footprint argument
preserves a previously restored compact inner circle (and the \(a\)-branch under
its own \(b\)-strip correction) when the source strip avoids it by a strict
margin. The strict margin holds for the projected footprint of the entire
reference branch or circle in the relevant slot.

\section{Differentiating the link splitting}\label{sec:splitting}

\subsection{A chart on an unchanged initial
interval}\label{a-chart-on-an-unchanged-initial-interval}

For the active link choose native first and translated
fundamental intervals on the upper leading segment. Their charts have
the form \((t,e)\mapsto(t,y_{\mathrm{top}}+e)\), with signed step
\(\lambda=-h\) or \(+h\). Pull them back by \(d_{\mathrm{lead}}\) model
iterates to obtain \(\theta_*\) inside the protected \(P\) neighborhood.
The entire closed intervals enter a local unstable seed box under
sufficiently many backward iterates. This holds uniformly on each
closed subarc because the iterates preserve order and both extreme
points tend to \(P\). Shrink the phase and energy collars to keep
positive margins. The actual return \(F_{j,V}\) coincides with the
baseline on the protected neighborhoods of \(P\) and \(Q\), for every
active \(V\) in the common small low-norm ball.

Translate the phase origin temporarily and take the \(+h\) direction. In the ideal chart, the
zero-active-profile return \(G_j\) tends to translation by \(h\). Build a smooth
coordinate change \(H_j\) on a neighborhood of \([0,h]\times[-a,a]\) that is the
identity on the collar of \(t=0\) and equals \(T_h G_j^{-1}\) on the collar of
\(t=h\), where \(T_h(t,e)=(t+h,e)\). One obtains it by interpolating the
difference from the identity using a fixed cutoff supported at the second
collar, with \(a\) and the collar widths shrunk first. Its determinant \(D_j\)
is positive and equals 1 on both collars. The map
\(K_j(t,e)=(t,\int_0^e D_j(t,s)\,ds)\) has the same pulled-back area form as
\(H_j\) and is identity on the collars. Then \(H_j K_j^{-1}\) is symplectic on a
slightly smaller common strip. Extend it into the next strip by
\(T_h(H_jK_j^{-1})G_j^{-1}\), which agrees identically on the overlap by the
collar choices. This constructs an exact time-energy chart for the return from the chosen fundamental strip into its image.
Writing \(\theta_j\) for the resulting symplectic embedding in the
return domain gives
\[
 F_{j,V}\theta_j(t,e)=\theta_j(t+\lambda,e)
\]
on the collared first strip for the whole active family. For \(-h\) the seam construction
runs in the other direction. The charts and their inverses converge in every
fixed norm to the ideal chart, because the interpolation, integral correction
and coordinate inversion are on fixed neighborhoods with uniform inverse bounds.

The unstable graph \(e=u_j(t)\) in the resulting chart is fixed during this
stage. The stable graph \(e=s_{j,V}(t)\) is obtained from the unchanged local
stable seed at \(Q\) by a fixed number \(L\) of inverse iterates, where \(L\) is chosen at the source on the two closed phase intervals
with collars. The images at \(L\) and \(L+1\) of the reference arc
lie inside the smaller \(Q\) seed box with a margin, by forward
convergence of the extreme points and order preservation. The margin
persists for the perturbed roots on the same fixed collars. An exact
scalar implementation uses forward iterates: in the local seed chart, subtract
the seed graph value from the second coordinate of \(F_{j,V}^{L}\theta_j(t,e)\). At the
source its \(e\)-derivative is nonzero, since the image of the tangent of the
initial branch is the tangent of the stable seed and the map is invertible.
There is a uniform lower bound by compactness. A fixed low-norm neighborhood of
zero active controls preserves this and puts the root in a fixed strip about the
source graph. Thus a scalar implicit equation gives the stable graph on the full
phase interval. Choose the protected boxes so the reference map and its inverse
are \(C^1\)-close to their hyperbolic linear parts in coordinates adapted to the
local stable and unstable directions. The zero-active-profile maps converge
there in every \(C^r\) norm, also after recentering at the continued saddles. On
smaller boxes the hyperbolic graph transform gives smooth local seeds and
convergence in each such norm \cite[Theorems 2.3--2.4]{HP}. The support
exclusion keeps these same seeds fixed as the active profile varies.

The splitting function \(M_{j}(V)=u_j-s_{j,V}\) is smooth periodic of period
\(h\). The two graph germs at the phase cut match because the chart conjugates
the return to translation and both branches are invariant. Its mean is zero. In
a primitive agreeing with \(e\,dt\) on the chart, a fundamental unstable segment
has integral \(S(z_u)-S(P)\) and a fundamental stable segment integral
\(S(z_s)-S(Q)\), where \(S\) is the action for one normalized return. This
follows by applying the action differential to a branch path from \(P\) or
\(Q\). The chart translation makes \(S(z_u)=S(z_s)\) on the connected strip.
Therefore the integral of the oriented graph difference is \(S(Q)-S(P)\), which
vanishes by the exact calibration. The reverse phase direction merely reverses
the sign of the integral relation. The local gauge change from \(Y\,dX\) is
exact and leaves fixed-point action differences unchanged.

\subsection{First variation of the stable
graph}\label{first-variation-of-the-stable-graph}

Let \(F=F_{j,V}\) and vary the active profile in a direction \(U\). In
normalized phase coordinates, write \(\partial_V F[U]\circ F^{-1}=X_K\), with
\(\iota_{X_K}\omega=dK\), \(\omega=de\wedge dt\) in the phase chart.
Differentiating \(F^{-L}\) applied to the fixed seed gives the Hamiltonian
\(-\sum_{\ell=1}^L K\circ F^\ell\) on the initial graph. Indeed a variation of
\(F^{-1}\) is \(- (F^{-1})_*X_K=X_{-K\circ F}\), and differentiating the \(L\)
factors gives the stated sum. A curve variation by \(X_H=(-H_e,H_t)\) changes a
graph at fixed \(t\) by \(H_t+s'(t)H_e=\partial_t[H(t,s(t))]\). Hence \[
 D M_j(V)[U](t)=\partial_t\sum_{\ell=1}^L K(F^\ell(t,s_{j,V}(t))).
\] A moving point of the seed only adds a tangent variation, which gives zero
graph increment in this formula.

For a physical reflection at the current curve, the post-relative Hamiltonian
variation is \(2\sin\varphi\) times the normal variation at the current impact.
This follows by the calculation of Section~\ref{sec:reflection}
applied at the current curve. A tangential displacement changes only the
first-order parametrization, hence the reference-normal direction contributes
\(n_E\cdot n_{\mathrm{current}}\) times its scalar profile. Composing the
physical flights and reflecting maps adds the transported Hamiltonians, and the
fixed affine scale divides them by \(J\). Since only the chosen active slots can
vary, the exact derivative is therefore the phase derivative of a finite sum of
terms \[
 a_{\ell i,j,V}(t)\,U_i(\sigma_{\ell i,j,V}(t)).
\] Here \(i\) labels a physical profile slot and \(\ell\) a normalized return
visit of the stable branch. The coordinate \(\sigma\) is the actual impact read
in the scaled inverse old stable-footpoint chart. The coefficient \(a\) is
\(2\sin\varphi(n_E\cdot n_{\mathrm{current}})\) divided by the fixed
\(2\sin\varphi_{\mathrm{old}}\) in the definition of that physical profile, with
the fixed support cutoff when present. The signed \(J\) cancels. At the ideal
model the uncut coefficient is 1. A later visit at a different phase to the same
physical slot uses the same function \(U_i\).

The same unrolled stable trajectory at \(t\) and \(t+\lambda\) starts one map iterate
apart, by exact translation on the chart. The two ends of the finite visit lists
can be placed in the protected neighborhoods where all active profile
evaluations are zero. Thus controlled visits at the cut match by shifting the
iterate labels, with equality of the projected impact maps, weights and their
germs. This establishes gluing of the actual terms as phase patches. The finite
sum may therefore be evaluated on either phase patch.

\subsection{Leading and folded visits}\label{leading-and-folded-visits}

For the active block \(H_{\widehat V'}R^3\) let \(I_2=[A,A+2h]\) be
the two-period interval on the straight leading segment. Parametrize
the primitive on a fixed large circle in \(X\). Fix \(0<\delta<h/4\),
take \(K=[A+\delta,A+2h-\delta]\), and put
\(K^+=[A+3\delta/4,A+2h-3\delta/4]\).
The fixed strip cutoff is supported inside \(I_2\) and is one on
\([A+\delta/2,A+2h-\delta/2]\).

At return index \(q\) in the sum above, start from
\(F_{j,V}^{q-1}\theta_j(t,s_{j,V}(t))\). In the reference model the
observed \(X\) is the first coordinate of \(F_*^q\theta_*(t,0)\).
Writing \(d=d_{\mathrm{lead}}\) locally, the potentially contributing
indices on the first phase interval are
\[
\begin{array}{ccl}
\text{branch}& q & \text{reference }X\\
a & d,\ d+1 & t,\ t-h\\
b & d,\ d+1 & t,\ t+h\\
b & d+4+r,\quad 0\le r\le 3
    & x_b+2h-(t-x_b)/2-rh/2 .
\end{array}
\]
On the translated phase patch the exact graph translation shifts the
starting orbit by one return. Thus its \(q\)-th sample is the
\((q+1)\)-st sample on the first patch, with the same weight and physical
coordinate. The unmatched end samples start in the protected boxes.
Every other reference sample has a positive margin to the support
of the strip cutoff, including any hit at a strip edge. The
convergence of the fixed finite itinerary and the impact coordinates
retains this margin on a late tail and common low ball.

The two successive leading visits per period have ideal coordinates \(t,t+h\) in
the increasing case (and \(t,t-h\) in the decreasing case). Define the actual
leading footpoint as a function \(\sigma^L(x)\) on the two-period interval: use
the first visit for its first phase sheet and the next visit for its other
sheet. At their shared cut, these two formulas agree on a neighborhood, because
advancing from one phase representative to the next applies exactly one map
return on the stable branch in the unchanged upstream chart. Their full germs,
including the action weights \(a^L(x)\), agree. Thus the two visits define one
smooth leading parametrization on both sheets. This is an unfolded interval of
physical control \(X\)-coordinates. The two sheets give distinct physical
coordinates in the same slot. For \(j\) large and all sufficiently small
controls in one fixed low norm, \(\sigma^L\) is strictly monotone, close to the
identity in \(C^1\). Its inverse is defined on a common collar of the primitive
support. The uncut leading coefficient \(a^L\) is at least \(1/2\) on the common
primitive collar on the same late tail and fixed low ball: at zero profiles it
tends uniformly to 1 by the weight calculation below, and uniform \(C^0\)
dependence allows shrinking the ball. The images of the support with its collar
lie in the fixed cutoff plateau by the same \(C^0\) margin. One physical
function thus determines both leading contributions.

For support verification, in the source the selected leading segment, and for
\(b\) the folded segment, are the only points of this link with \(X\) in the
active strip. The remaining source hits in the fixed finite iterates stay a
positive distance from the smaller active support. The finitely many model
rotations of the one identity word and the previously fixed functions are
controlled on fixed compacts. The diagonal family retains these margins, by
convergence of the return and of the normalized physical impact readouts. Where
the endpoints of the leading and folded segments are used, work on collars in
their smooth prepared neighborhoods. The invariant relation at the cut makes the
definition smooth across phase sheets. At points whose shifted supports enter no
visit, extend the control contribution by zero through its fixed smooth collars.

The normalized physical impact readout is the inverse old stable-footpoint
coordinate of the actual scalar root in the reflection calculation. At a
reference path its difference from the normalized old output horizontal
coordinate is \(O_{C^r}(\tau/\mu+\tau)\) for bounded fixed-order profiles: the
\(\tau/\mu\) term is the old footpoint's dependence on the raw transverse
height, and the \(\tau\) term is the exact rescaled impact displacement. The
positive reflection weight is read from the incoming direction, the current
normal, and the fixed old weight, all in uniformly transverse charts. At zero
profiles its ratio to the old weight tends to 1 in every fixed norm. The same
coefficient and root formulas give uniform tame bounds on a fixed low-norm
neighborhood as in the reflection calculation. The stage maps, actual impact
readouts and weights are then evaluated along the finitely many normalized
iterates of the stable graph, computed by the scalar root above, and along the
fixed suffix. The model prefixes and all frozen charts are uniformly smooth in
the row. The number of variable and normalized stages in these evaluations is
fixed, even though the prefix lengths vary.

For the tame calculus here, products and finite compositions are estimated by
the chain rule with interpolation between a fixed low seminorm and the highest
needed seminorm. A derivative in each profile direction contributes that
direction once. For scalar roots or one-dimensional inverse parametrizations,
differentiation produces a uniform reciprocal of the nonvanishing root derivative
times the lower terms, allowing the same jet count as in the reflection
calculation. More explicitly, for composition the total derivative order
allocated to derivatives of the inside map beyond its first derivative, plus
derivatives of the outside function beyond its first derivative, is at most the
total order minus one for state derivatives. New parameter directions increase
this budget by at most one each when their ordinary derivative orders are
included. The same accounting in \(D(G^{-1})=1/(G'\circ G^{-1})\) or in the
differentiated implicit identity, with base low bounds and a fixed positive
denominator, supplies a loss independent of the output order. Iterating this
over a fixed number of operations gives \(C^r\) tame estimates with input
\(C^{r+d_q}\) for each fixed number \(q\) of parameter derivatives. It applies
to the stable graph, the coefficient functions, the lifted leading map and its
inverse, and the folded maps below. The coincidence of the germs at a seam
allows all estimates on a fixed finite set of overlapping patches.

\subsection{Exact leading alignment and the folded
operator}\label{exact-leading-alignment-and-the-folded-operator}

Choose a fixed nonnegative smooth \(\rho\), supported in the interior
of \(K\), with \(\rho(t)+\rho(t+\lambda)=1\) on the first phase
interval. Such a function is obtained by normalizing a nonnegative
bump under periodization by \(h\), since the interior of \(K\)
contains representatives of every phase. The actual leading map on
the collared \(I_2\) has inverse on \(K^+\), and the images of \(K^+\)
under \(\sigma^L\) lie in the cutoff-one interval by the strict margins
and \(C^1\) closeness to the identity. For a smooth zero-mean forcing \(g\) of period \(h\), define its supported
primitive \[
 V_g(x)=\int_{-\infty}^{x}\rho(v)g(v)\,dv.
\] The full integral is zero, so this primitive is supported in the same closed
interior interval. It gains one derivative in all the fixed norms. Use the
actual direction \[
 (R_Vg)(\sigma^L(x))=V_g(x)/a^L(x),
\] extended by zero through its collars on the large input circle. Nearby
leading images of this support remain inside the region where the fixed profile
cutoff equals 1. Thus this is a linear-in-\(g\), smooth tame family of bona fide
physical primitive directions. The two leading action contributions sum to
\(V_g(t)+V_g(t\pm h)\), and taking the phase derivative produces \(g(t)\)
exactly. This yields a right inverse on the \(a\)-link, which has no other
controlled visits.

For the \(b\)-link, four further visits occur on the folded segment. At the
source, for \(t\in[x_b,x_b+h]\), their coordinates are \[
 z_\ell^0(t)=(3x_b+\ell h-t)/2,\qquad 1\le\ell\le4,
\] by the lower time-energy chart in Section~\ref{sec:model} and
its four-period folding \cite[Definition 4.4 and Lemma 4.5]{BT}. The ratios of
folded to leading weights are 1. In the actual family use \[
 z_\ell(t)=(\sigma^L)^{-1}(\sigma_\ell(t)),\qquad
 c_\ell(t)=a_\ell(t)/a^L(z_\ell(t)),
\] on the needed enlarged patches, extending through zero-support collars. The
actual \(z_\ell,c_\ell\) are uniformly tame as just shown, and are arbitrarily
\(C^2\)-close to \(z_\ell^0,1\) by first taking the diagonal tail and then
shrinking a fixed low-norm ball of profiles. The \(b\)-link derivative on
\(R_Vg\) is \[
 D M_j(V)R_Vg
 =g+\sum_{\ell=1}^4 c_\ell\rho(z_\ell)z_\ell' g(z_\ell)
     +\sum_{\ell=1}^4 c_\ell' V_g(z_\ell).
\] There is no extra primitive-weight term at the source, where the coefficients
are constant. Each actual term comes from the same physical primitive at its own
reparameterized impact, so the formula keeps the repeated-visit constraint. The
operator on the left is periodic and maps into zero mean. The principal
perturbation contains contracted phase arguments, clustered in pairs that differ
by \(h\) at the source. The quantitative tame inverse of this operator is the
input to smooth exact repair.

\section{Smooth restoration of the links}\label{sec:restoration}

For the ambient-map restoration in \cite[\S5]{BT}, the authors' corrigendum uses
a compactness argument with a loss of one derivative \cite{BTcorr}. Here the
splitting equations for the billiard profiles use fixed-loss derivative bounds
and tame local surjectivity.

\subsection{Inverting the folded
contribution}\label{inverting-the-folded-contribution}

Let \(E_h\) be the zero-mean smooth functions of period \(h\). Use
\(|g|_1=\|g'\|_\infty\). This is a norm on the zero-mean \(C^1\) completion. Let
\(T=I-DM(V)R_V\) for the \(b\)-link and write \[
 Tg=\sum_{\ell=1}^4 a_\ell(t)g(z_\ell(t))
      +\sum_{\ell=1}^4 b_\ell(t)V_g(z_\ell(t)),\quad
 a_\ell=-c_\ell\rho(z_\ell)z_\ell',\quad b_\ell=-c_\ell'.
\] The full \(T\) is periodic and preserves zero mean, as follows from the
actual splitting map. The summands may be interpreted on phase patches and need
not be individually periodic. At the source \(z_\ell'= -1/2\), \(c_\ell=1\),
\(z_{\ell+2}^0=z_\ell^0+h\) for \(\ell=1,2\), and \(a_\ell^0+a_{\ell+2}^0=1/2\).
In particular \(\sup_t\sum |a_\ell^0|=1\).

Take the \(C^2\) error in \(z_\ell,c_\ell\) at most \(\delta\). In the
derivative of the principal sum the terms involving \(g'\) have bound
\((1/2+C\delta)|g|_1\). For each pair write its coefficient-derivative
contribution as \[
 (a_\ell'+a_{\ell+2}')\,g(z_\ell^0)
 + a_\ell'[g(z_\ell)-g(z_\ell^0)]
 + a_{\ell+2}'[g(z_{\ell+2})-g(z_{\ell+2}^0)].
\] The coefficient sum is \(O(\delta)\), since the source pair sum is constant.
The differences are \(O(\delta)|g|_1\), and \(\|g\|_\infty\le h|g|_1\). The
derivative of the primitive-order sum is \(O(\delta)|g|_1\) because
\(b_\ell,b_\ell'=O(\delta)\) and \(\|V_g\|_{C^1}\le C|g|_1\). Hence
\(|Tg|_1\le q_0|g|_1\) with a fixed \(q_0<1\) on a common sufficiently small
low-norm neighborhood in a sufficiently late tail. The \(C^1\) inverse is the
Neumann series \(S=(I-T)^{-1}\).

\subsection{Higher regularity with fixed
loss}\label{higher-regularity-with-fixed-loss}

There are uniform constants \(M,\kappa\), with \(M\kappa<1\) and \(\kappa<1\),
bounding \(\sum_\ell |a_\ell|\) and \(\max_\ell|z_\ell'|\), respectively. For
\(r\ge2\) the top term in the principal sum has bound
\(M\kappa^r\|g^{(r)}\|_\infty\). Every other term there involves \(g^{(p)}\)
with \(p<r\). The primitive-order terms also have this property because
\(V_g'=\rho g\).

Here is the quantitative lower-term estimate, with arbitrary \(\eta>0\).
Interpolate the derivatives of \(a_\ell,b_\ell,z_\ell\) between \(C^1\) and
\(C^r\), and those of \(g\) between \(C^1\) and \(C^r\). A term
\(a_\ell^{(l)}g^{(p)}(z_\ell)\prod_{s=1}^p z_\ell^{(m_s)}\), where
\(l+\sum m_s=r\) and \(m_s\ge1\), has total interpolation weight at most \(r-1\)
before division by \(r-1\). The weights are \((l-1)_+\), \((p-1)_+\) and
\(m_s-1\), with the case \(p=0\) interpreted without products. In the
\(g\)-factor its high exponent is strictly less than one if \(p<r\). For the
primitive term use \(V_g\) between \(C^2\) and \(C^{r+1}\), whose bounds are
controlled by those of \(g\). This reduces the weight in that factor further.
The finite-interval interpolation inequality from the reflection proof and
Young's inequality bound the sum of non-top terms by \[
 \eta\|g^{(r)}\|_\infty+
 C_{r,\eta}(1+\mathcal A_r)|g|_1,
 \qquad \mathcal A_r=\max_\ell\|(a_\ell,b_\ell,z_\ell)\|_{C^r}.
\] Here one may first use the full \(C^r\) norm of \(g\), then use its
equivalence to the highest derivative seminorm on zero-mean periodic functions
and choose the Young coefficient accordingly. The low coefficient bounds are
uniform, including for the lifts of \(z_\ell\) on fixed phase patches. The
relevant derivatives of the primitive are bounded on the fixed enlarged
projection interval. The coefficients have
\(\mathcal A_r\le C_r(1+\|V\|_{C^{r+d}})\) for a fixed loss \(d\), by the link
and readout estimates (including the one derivative in \(b_\ell\)).

Choose \(\eta\) so \(M\kappa^r+\eta\le q_1<1\) for all \(r\ge2\). Keeping the
lower estimate and iterating the equation for the Neumann partial sums yields
uniform bounds in each \(C^r\) norm. Compact embedding from \(C^{r+1}\) and the
unique \(C^1\) limit then show the limit is smooth and give \[
 \|Sg\|_{C^r}\le C_r\left(\|g\|_{C^r}
             +(1+\|V\|_{C^{r+d}})\|g\|_{C^1}\right).
\] All constants are uniform in the row. The shift is independent of the output
order \(r\).

The same expression for \(T\) and the finite composition estimates show
\(T(V)g\) is a smooth tame map in \((V,g)\) with a finite shift for each
parameter derivative. Resolvent differences satisfy
\(S(V_2)-S(V_1)=S(V_2)(T(V_2)-T(V_1))S(V_1)\). The inverse bounds and this
identity first give continuity in the smooth topologies and then justify
differentiating to obtain \(D S(V)[U]g=S(V)(D T(V)[U])S(V)g\). Iterating this
formula gives each higher derivative by finite products of the same tame maps,
with total shift depending on the derivative order in \(V\) and independent of
the output order. Thus \(R_V S(V)\) is the required uniform smooth tame family
of linear right inverses. The \(a\)-link uses \(R_V\) without the inverse
correction.

\subsection{The splitting equation on a space of vanishing
sequences}\label{the-splitting-equation-on-a-space-of-vanishing-sequences}

The input space for this stage is \(C^\infty(\mathbb T_X;c_0)\), where
\(\mathbb T_X\) is the fixed large circle used for the free primitive and
\(c_0\) denotes scalar real sequences vanishing at infinity. The output is the
zero-mean subspace of \(C^\infty(\mathbb R/h\mathbb Z;c_0)\). Equivalent
seminorms are \(\sup_j\|\cdot\|_{C^r}\), with component sequences vanishing in
each \(C^r\) norm. For the latter characterization, compactness of the phase
circle and the definition of the \(c_0\) norm give uniform tail convergence of
each Banach-valued derivative, and the reverse implication follows by uniform
Taylor estimates in the next derivative seminorm.

These spaces are tame Fréchet spaces. For a Banach value space \(B=c_0\), use a
smooth frequency partition \(\Delta_m\) on scales \(e^m\) and spectral cutoff
\(\widetilde\Delta_m=1\) on each piece's multiplier support. Map \(f\) into
\((\Delta_m f)_m\) and reconstruct by \(\sum_m\widetilde\Delta_m f_m\). With
\(B'=C^0(\mathbb T;B)\), the sequence norms \(\sum_m e^{mn}\|f_m\|_{B'}\) use
the \(\Sigma(B')\) exponential grading. Passing between its sum and supremum
versions costs one additive grading shift. Smooth-band convolution kernel bounds
and integration by parts give \[
 \sum_me^{mn}\|\Delta_m f\|_{B'}\le C_n\|f\|_{C^{n+2}},
 \qquad
 \left\|\sum_m\widetilde\Delta_m f_m\right\|_{C^n}
       \le C_n\sum_m e^{mn}\|f_m\|_{B'}.
\] The kernel bounds hold for Banach values by the integral triangle inequality.
Reconstruction on the range is the identity by the partition. Taking zero mean
is a continuous tame projection. This gives the tame direct summands required by
the theorem.

Fix one stage with the same protected phase and seed family as in
Section~\ref{sec:splitting}. On a late tail \(j\ge j_0\) choose one
integer \(b\) large enough for the scalar root estimates and the
\(C^2\) folded bounds. The common domain is the open ball
\[
 {\cal U}=\{\mathbf V=(V_j):\ %
      \sup_{j\ge j_0}\|V_j\|_{C^b}<\varepsilon\}
\]
in the input tame space. Choose \(\varepsilon\) within the uniform row
radius. Write \({\cal R}_j(V)\) for \(R_V\) on the \(a\)-branch and
for \(R_V S(V)\) on the \(b\)-branch, with the data of row \(j\).
The component maps are
\[
 {\cal M}(\mathbf V)_j=M_j(V_j),\qquad
 (\Psi(\mathbf V)\mathbf g)_j={\cal R}_j(V_j)g_j .
\]
The periodic action identity gives zero mean for every \(M_j(V_j)\).
At zero active profile the protected charts, local seeds and fixed
finite propagation converge to the same reference arc. Thus the
residuals \(M_j(0)\) tend to zero in each fixed norm. For a fixed
output order \(r\), the uniform first derivative and inverse
estimates bound the two increments by
\[
\begin{split}
 \|M_j(V_j)-M_j(0)\|_{C^r}
   &\le C_r\big[\|V_j\|_{C^{r+d}}
       +(1+\|V_j\|_{C^{r+d}})\|V_j\|_{C^b}\big],\\
 \|{\cal R}_j(V_j)g_j\|_{C^r}
   &\le C_r\big[\|g_j\|_{C^{r+d}}
       +(1+\|V_j\|_{C^{r+d}})\|g_j\|_{C^b}\big].
\end{split}
\]
Here a fixed shift \(d\) can be increased to cover the orders below
\(b\). The first inequality follows by integrating the derivative
along the segment from zero to \(V_j\). The two bounds show that
the outputs vanish in every fixed tail norm. For higher
derivatives the corresponding multilinear bounds keep one factor
from each direction. Taylor remainders estimated with the next
derivative are uniform in the row in a finitely higher seminorm.
This makes the componentwise derivatives actual Fréchet derivatives
and gives smooth tame evaluation of \({\cal M}\) and \(\Psi\).
The row identity gives \(D{\cal M}(\mathbf V)\Psi(\mathbf V)=I\) on the
output space throughout \({\cal U}\). Proposition 51.19 of \cite{KM}
now gives local surjectivity at zero between these tame Fréchet spaces.

Let \(r_j=M_j(0)\). Form target sequences equal to \(r_j\) up to a variable
cutoff and zero thereafter. They converge to \(r=(r_j)_j\) in the same smooth
zero-mean sequence space. For a sufficiently late cutoff, local surjectivity
produces a preimage \(V\) whose entries are smooth primitives and tend to zero
in each fixed norm, and which solves \(M_j(V_j)=0\) for every row in the tail.
The exact invariant graphs then coincide on a full fundamental interval, so
their stable and unstable extensions supply the whole heteroclinic branch.

Carry this out for the \(a\)- then \(b\)-branch of each inner circle, and last
for the outer circle, using finitely many fixed slot groups. At each step the
charts and seeds are fixed for that active family, and their convergence, the
actual zero-mean splitting and the tame right inverse used in Proposition 51.19
all refer to the same diagonal rows. Upon freezing the vanishing solutions, the
inverse chart and the actual finite forward and backward propagations on the
fixed covering patches converge in the physical model coordinates. At every
stage the calibrated actions and protected local end behavior remain exact. The
current suffix supports avoid all previously restored inner circles and, for a
\(b\)-repair, its companion \(a\)-branch. This is a physical itinerary property
from the strict projected-support margins, so later active controls preserve the
earlier branches exactly. All frozen corrections tend to zero in every fixed
norm. After taking a common tail, every row has all four exact smooth bi-links
with the source ordering.

\section{Positive entropy of the collision maps}\label{sec:entropy}

Let \(F_j\) be the normalized quotient return of one actual table after the
successive smooth repairs. Write \(Q_j\) for the first landing chart and \(n_j\)
for the number of physical quotient reflections per full word. The tables remain
strictly positively curved and smooth. The actual maps
\(F_j=Q_j^{-1}\overline T_j^{\,n_j}Q_j\) converge smoothly to the prepared model
on a fixed neighborhood of its outer circle and enclosed disc. The repairs give
four smooth invariant bi-links with the model ordering. Their local portions at
the eight saddles converge in \(C^2\) by the local stable and unstable manifold
theorem for the nearby smooth maps. The remaining compact portions are obtained
by a fixed finite number of iterates and converge in \(C^2\) as well. The two
halves at \(P_i\) belong to the same local unstable manifold, and the two halves
at \(Q_i\) belong to the same local stable manifold. Thus each whole circle is
smoothly embedded, including across its two endpoints. Inner and outer ordering
follows from the strict separations of the four reference circles and their
fixed small tubular neighborhoods.

To apply stability of the reference island, extend the comparison on its
original closed disc. On the fixed larger rectangle from
Section~\ref{sec:model}, the map
\(\Psi_j^{\mathrm{loc}}=F_*^{-1}\circ F_j\) is symplectic and converges smoothly
to the identity. Write it as \((x,y)\mapsto(X,Y)\). The derivative
\(\partial Y/\partial y\) is near one, and the mixed coordinates \((x,Y)\) cover
a fixed smaller rectangle with a collar containing the prepared source disc. The
closed one-form \(y\,dx+X\,dY\) has primitive \(xY+\eta_j(x,Y)\), where
\(\eta_j\to0\) in every finite norm after normalization at a fixed base point.

The closed region inside the prepared outer circle lies compactly in
\(\Phi(\operatorname{int}D_{\mathrm{BT}})\). Take a fixed cutoff \(\chi\)
supported inside this latter domain and equal to one near the closed region with
a fixed margin. Extend \(\chi\eta_j\) by zero in mixed coordinates. The
equations \[
 y=Y+\partial_x(\chi\eta_j)(x,Y),\qquad
 X=x+\partial_Y(\chi\eta_j)(x,Y)
\] give a global symplectic diffeomorphism \(\Psi_j\). For fixed \(x\), the
first equation is strictly increasing and onto as a function of \(Y\). For fixed
\(Y\), the second is strictly increasing and onto as a function of \(x\). Both
follow from \(C^2\) smallness and compact support. Whenever the correction
differs from the identity, \((x,Y)\) belongs to the compact interior support and
\((x,y)=(x,Y)+o(1)\), so the correction is the identity near and outside the
boundary of the prepared source disc. It agrees with \(\Psi_j^{\mathrm{loc}}\)
near the four actual circles and the closed outer region for all late \(j\),
using the fixed margin and the \(C^2\) convergence of those curves.

Then \(G_j=F_*\circ\Psi_j\) agrees with the physical normalized return on that
neighborhood. Conjugating by \(\Phi^{-1}\) gives a conservative diffeomorphism
preserving \(D_{\mathrm{BT}}\) and \(C^2\)-close there to the source map. Its
four bi-links are \(C^2\)-close to the corresponding marked circles, with the
same continued saddles. By \cite[Proposition 2.1 and Corollary E]{BT} these
links bound an invariant stochastic island. In the prepared coordinates let
\(I_j\) be this open island, of positive area and positive maximal Lyapunov
exponent throughout. The comparison agrees with \(F_j\) on a neighborhood of
\(I_j\), which gives invariance and the same exponents for \(F_j\) there.

The comparison map is identity outside a compact subset of the plane. It
therefore defines a smooth area-preserving diffeomorphism of a sufficiently
large flat torus. The probability given by normalized area on \(I_j\) is
invariant and absolutely continuous with respect to ambient area. Pesin's
entropy formula for this probability \cite{Mane} expresses its entropy as the
integral of the sum of positive Lyapunov exponents with multiplicity. The
integral is strictly positive, hence the restriction of \(F_j\) to \(I_j\) has
positive entropy for normalized area.

The island \(Q_j(I_j)\) is invariant under \(\overline T_j^{\,n_j}\) and has
positive quotient Liouville area, since the affine canonical Jacobian is the
nonzero \(J_j\). The normalized restriction of quotient Liouville measure is
carried back by \(Q_j\) to normalized model area. To see positivity for the
entire quotient probability, choose a finite positive-entropy partition for the
normalized restriction, and extend it by one atom for the invariant complement.
If the island has probability \(a>0\), the entropy of each finite joined
partition is \(a\) times the corresponding entropy on the island plus the
entropy of the binary island/complement partition. Dividing by the time and
letting it tend to infinity gives strictly positive full probability entropy for
\(\overline T_j^{\,n_j}\). For any finite partition the \(m\)-step join under
the \(n_j\)-th iterate is coarser than the \(mn_j\)-step join under the original
map. Consequently positive iterate entropy implies positive entropy for
\(\overline T_j\). The normalized quotient area is the pushforward of the normalized
full ray area. The antipodal quotient has two sheets and preserves
area on each, so the degree of the cover cancels after normalization.
Pulling quotient partitions back to the full oriented-line
annulus preserves their joined probabilities and gives the same lower bound for
\(T_j\) with the normalized full Liouville probability.

In the collision chart let \(s\) be the arclength of the \textbf{actual}
boundary and \(\varphi\in(0,\pi)\) the angle of the outgoing direction from the
positive tangent. The tangential momentum is \(p=\cos\varphi\), unchanged by
reflection at the collision. The positive measure of the two-form
\(dp\wedge ds\) is \(\sin\varphi\,ds\,d\varphi\), and its total mass is
\(2|\partial D_j|\). The conjugacy from incoming lines to outgoing collision
states carries \(T_j\) to the physical specular collision map \(B_j\), with
precisely the probability in the target. Thus \(h_{\mu_j}(B_j)>0\) for all
sufficiently late rows.

Finally write \(\gamma_E(\theta)=(2\cos\theta,\sqrt3\sin\theta)\) on
\(S^1=\mathbb R/(2\pi\mathbb Z)\). Each boundary is
\(\gamma_E(\theta)+f_j(s_E(\theta))n_E(\theta)\), where \(s_E\) is fixed
reference arclength and \(f_j\) is the sum of the physical normal profiles. The
designated \(k_j\) suppresses all row-dependent finite coefficients through
increasing orders. The actual correction profiles vanish in every fixed norm in
the solved sequence. It follows that \(f_j\to0\) in every fixed \(C^r\) norm,
and hence the smooth embeddings converge to \(\gamma_E\) in \(C^\infty\).
Discard the finitely many initial rows and renumber. This proves
Theorem~\ref{thm:main}.


\subsection*{AI assistance}
AI tools assisted in developing and drafting the mathematical argument,
checking its natural-language dependencies, and reviewing the exposition.

\begin{thebibliography}{11}
\bibitem{BergerICM}
P.~Berger, \emph{Wild dynamics on manifolds}, Proceedings of the
International Congress of Mathematicians 2026, Vol.~5, 24--41.
\href{https://doi.org/10.1137/25M1810787}{doi:10.1137/25M1810787}.

\bibitem{BT}
P.~Berger and D.~Turaev, \emph{On Herman's positive entropy conjecture},
Adv. Math. \textbf{349} (2019), 1234--1288.
\href{https://doi.org/10.1016/j.aim.2019.04.002}{doi:10.1016/j.aim.2019.04.002}.

\bibitem{BTcorr}
P.~Berger and D.~Turaev, \emph{Corrigendum to On Herman's Positive
Entropy Conjecture}, appended to \href{https://arxiv.org/pdf/1704.02473v2#page=53}
{arXiv:1704.02473v2}, 2026.

\bibitem{Bessa}
M.~Bessa, G.~Del Magno, J.~Lopes Dias, J.~P. Gaiv\~ao and M.~J. Torres,
\emph{Billiards in generic convex bodies have positive topological
entropy}, Adv. Math. \textbf{442} (2024), 109592.
\href{https://doi.org/10.1016/j.aim.2024.109592}{doi:10.1016/j.aim.2024.109592}.

\bibitem{DCRR}
J.~Damasceno, M.~J. Dias Carneiro and R.~Ram\'irez-Ros,
\emph{The billiard inside an ellipse deformed by the curvature flow},
Proc. Amer. Math. Soc. \textbf{145} (2017), 705--719.
\href{https://doi.org/10.1090/proc/13351}{doi:10.1090/proc/13351}.

\bibitem{DRR}
A.~Delshams and R.~Ram\'irez-Ros, \emph{Poincar\'e--Melnikov--Arnold
method for analytic planar maps}, Nonlinearity \textbf{9} (1996), 1--26.
\href{https://doi.org/10.1088/0951-7715/9/1/001}{doi:10.1088/0951-7715/9/1/001}.

\bibitem{Glutsyuk}
A.~Glutsyuk, \emph{Density of thin film billiard reflection pseudogroup
in Hamiltonian symplectomorphism pseudogroup},
\href{https://arxiv.org/abs/2005.02657v5}{arXiv:2005.02657v5}.

\bibitem{HP}
M.~W. Hirsch and C.~C. Pugh, \emph{Stable manifolds and hyperbolic
sets}, Global Analysis, Proc. Sympos. Pure Math. \textbf{14} (1970),
133--163. \href{https://escholarship.org/uc/item/2rr3r1pz}{Primary text}.

\bibitem{KM}
A.~Kriegl and P.~W. Michor, \emph{The Convenient Setting of Global
Analysis}, Mathematical Surveys and Monographs \textbf{53}, AMS, 1997.
\href{https://www.mat.univie.ac.at/~michor/apbookh-ams.pdf}{Authors' edition}.

\bibitem{Mane}
R.~Ma\~n\'e, \emph{A proof of Pesin's formula}, Ergodic Theory
Dynam. Systems \textbf{1} (1981), 95--102.
\href{https://doi.org/10.1017/S0143385700001188}{doi:10.1017/S0143385700001188}.

\bibitem{Turaev}
D.~Turaev, \emph{Polynomial approximations of symplectic dynamics and
richness of chaos in non-hyperbolic area-preserving maps},
Nonlinearity \textbf{16} (2003), no.~1, 123--135.
\href{https://doi.org/10.1088/0951-7715/16/1/308}{doi:10.1088/0951-7715/16/1/308}.
\href{https://doi.org/10.20347/WIAS.PREPRINT.722}{WIAS Preprint 722}.
\end{thebibliography}
\end{document}
